\documentclass[10pt,twocolumn,letterpaper]{article}

% ============================================================
% Learning to Individuate --- Abstract + Introduction draft
% Self-contained conference-style preprint (compiles with pdflatex).
% Swap in cvpr.sty for the camera-ready.
% ============================================================
\usepackage[utf8]{inputenc}
\usepackage[T1]{fontenc}
\usepackage[margin=0.75in]{geometry}
\usepackage{times}
\usepackage{amsmath,amssymb}
\usepackage{graphicx}
\usepackage{booktabs}
\usepackage{xcolor}
\usepackage{tikz}
\usetikzlibrary{positioning}
\usepackage{microtype}
\usepackage[hidelinks]{hyperref}

\newcommand{\U}{\mathcal{U}}
\newcommand{\method}{\textsc{IndividuateNet}}
\newcommand{\tbd}{\textcolor{red}{\footnotesize TBD}}

\title{\textbf{Learning to Individuate: Information-Maximising Caption\\
Generation for Zero-Shot Face Enrolment}}
\author{Anonymous Author(s)\\ \textit{Paper ID ----}}
\date{}

\begin{document}
\maketitle

% ------------------------------------------------------------
\begin{abstract}
Can a face be enrolled from words alone, and can we say \emph{in advance} which words will
work? Recent work shows that a dual-encoder aligning face crops with natural-language
descriptions can verify \emph{unseen} identities from text under a strict open-set protocol,
and that its robustness to appearance change is governed almost entirely by a single data-side
property: whether captions \emph{individuate} --- describe the person rather than mutable
attributes. Yet individuation has remained an emergent, undefined property, supplied
heuristically by prompting an off-the-shelf vision--language model (VLM) on each identity's own
images, and measurable only \emph{after} a full training-and-evaluation cycle. We make
individuation a first-class, optimisable quantity. First, we define an information-theoretic
measure of individuation, $\U(T)=I(T;Y)-I(T;A\mid Y)$ --- the identity information a caption
carries net of the mutable-attribute information it leaks --- and show it is estimable directly
from an (image, caption, identity) corpus, \emph{without} training a biometric matcher. Second,
a short analysis links $\U$ to the empirically observed linear invariance law, predicting the
rate at which the attribute-invariance gap closes as captions individuate. Third, we introduce
\method, a caption generator trained to \emph{maximise} identity-discriminative information
while suppressing mutable-attribute content, which produces individuating descriptions
\emph{without} an external VLM at deployment. On the image-abundant regime (VGGFace2), \method{}
lifts zero-shot enrolment from chance to a strong operating point without VLM prompting, matches
or exceeds VLM-authored captions on CelebA, and its predicted individuation tracks measured
enrolment across three datasets spanning sparse/abundant imagery and attribute/individuating
captions. The takeaway reframes robust face--language matching as \emph{designing the
supervision} rather than engineering an invariance head: what the caption says can be measured,
predicted, and now generated.
\end{abstract}

% ------------------------------------------------------------
\section{Introduction}

Consider an investigator who must find a person from a witness's account when no photograph
exists. Conventional face recognition cannot begin here: it enrols an identity by storing one or
more reference \emph{images} and later comparing a probe against that
gallery~\cite{arcface,facenet}. This image-centric design is powerful but rigid --- enrolment
requires a photograph, the stored template is an opaque vector, and the enrolled appearance is
frozen at capture time --- yet many consequential settings supply only a \emph{verbal}
description: a witness statement, a case file, a watch-list attribute sheet. A recent line of
work shows that such a description can nonetheless serve as a first-class biometric enrolment: a
dual-encoder that aligns face crops with structured descriptions verifies a \emph{previously
unseen} identity from text alone, under a strictly open-set protocol, at operating points
approaching image--image recognisers~\cite{semanticbio}.

That same work asks \emph{why} a text-enrolled identity survives changes in appearance, and
returns a striking answer: robustness is governed not by the model but by the words. The
decisive property is whether a caption \emph{individuates}. Contrast two descriptions of one
person: ``brown hair, glasses, a beard'' names a \emph{type} that thousands share and that
changes with a haircut, whereas ``deep-set eyes set unusually wide, an aquiline nose with a
slight leftward deviation'' names \emph{that individual} through traits that persist across
appearance change. As captions move from the former, attribute-enumerating style toward the
latter, individuating style, the attribute-invariance gap --- the error increase between
consistent-look and changed-look verification of the same identity --- collapses near-linearly,
and it does so \emph{regardless} of any embedding-side disentangler, which is inert in isolation.
What the caption \emph{says} matters far more than how the embedding is regularised.

This finding is as unfinished as it is striking, because individuation is named as the causal
variable but never made concrete. It is not \emph{defined}: no quantity separates an
individuating description from an attribute-enumerating one. It is not \emph{measured}: it can be
read off only \emph{post hoc}, through the invariance gap, after a complete
training-and-evaluation cycle. It is not \emph{explained}: the near-perfect linearity of the
invariance law is reported as an empirical regularity, not derived. And it is not
\emph{generated} in any principled way --- individuating captions are obtained by prompting an
off-the-shelf vision--language model (VLM) on each identity's own images, an uncontrolled,
unoptimised, closed dependency that is the acknowledged central limitation of the approach and
that fails outright when per-identity imagery is scarce. If individuation is genuinely the lever,
it should be possible to \emph{define} it, \emph{predict} outcomes from it, and \emph{optimise}
for it directly.

This paper does exactly that, treating individuation not as an emergent by-product of prompting
but as a measurable, predictable, and \emph{learnable} property of the language supervision. Let
$I$ be a face crop, $T$ a description, $Y$ the identity, and $A$ a vector of \emph{mutable}
attributes (hair, makeup, accessories, expression). We define the \emph{individuation functional}
\begin{equation}
  \U(T) \;=\; \underbrace{I(T;Y)}_{\text{identity-bearing}} \;-\;
              \underbrace{I(T;A\mid Y)}_{\text{mutable-attribute leakage}},
  \label{eq:U}
\end{equation}
the identity information a caption carries \emph{net of} the transient-appearance information it
leaks: an attribute-enumerating caption scores high on $I(T;A)$ but low on the identity residual
$I(T;Y)$, and an individuating caption inverts this. Both terms are estimable from an
(image, caption, identity) corpus with standard neural mutual-information estimators, so $\U$
acts as a cheap, \emph{training-free} diagnostic computed \emph{before} any biometric matcher is
built. A short analysis then connects $\U$ to the observed invariance law: under a nested
caption-dose schedule the expected changed-look score shift is affine in the individuating
fraction $p$, yielding $\mathrm{gap}(p)=g_0-\kappa p$ with
$\kappa\propto(\U_{\text{rich}}-\U_{\text{disc}})$ --- an account of both the slope and the
striking linearity previously reported only as an empirical fit. Finally, because $\U$ is
differentiable through its estimators, it is not merely a diagnostic but an \emph{objective}: we
train \method, an adapter-tuned caption generator that maximises identity-discriminative
information while penalising mutable-attribute content and constraining outputs to fluent,
human-plausible language. Having internalised \emph{what to say} in order to individuate,
\method{} needs no per-identity image gallery and no external VLM at deployment.

Our claim is deliberately narrow and testable. We do not seek to beat dedicated image--image
recognisers on accuracy; we study a capability they do not possess --- text-only enrolment ---
and show that its central, previously fragile ingredient can be measured, predicted, and
generated. Concretely, we ask: (i)~does $\U$ separate individuating from attribute-enumerating
captions and \emph{predict} downstream enrolment across data regimes, without biometric training?
(ii)~does the $\U$-derived law quantitatively match the measured rate at which the invariance gap
closes? and (iii)~can a generator trained on $\U$ replace VLM-prompted captions, recovering
enrolment in the image-abundant regime \emph{without} the external model?

We answer all three affirmatively. Estimated $\U$ tracks measured zero-shot AUC and the
invariance gap across three datasets that bracket a $2\times2$ of \{sparse, abundant\} imagery
$\times$ \{attribute, individuating\} captions, reproducing the known regime map with no
per-configuration training. The $\U$-predicted slope agrees with the measured caption-dose
response. And \method{} lifts zero-shot enrolment on the image-abundant regime from chance to a
strong operating point \emph{without} any VLM at deployment, matching or exceeding VLM-authored
captions on the dense-attribute regime, while an attribute-collision stress test and a human
fluency evaluation confirm the gains reflect genuine individuation rather than attribute
enumeration or reward-hacked verbosity. The broader lesson is a shift of emphasis: robust
face--language matching is best pursued by \emph{designing the supervision} --- deciding what the
caption should say --- rather than by engineering an invariance head into the embedding.

\vspace{2pt}
\noindent\textbf{Contributions.}
\begin{itemize}\itemsep2pt
  \item \textbf{A measure of individuation.} We define $\U$
        (Eq.~\ref{eq:U}), the first information-theoretic, \emph{training-free} quantification of
        how much a caption individuates an identity, and validate it against controlled corpora.
  \item \textbf{A predictor, and an explanation of the invariance law.} We show $\U$ predicts the
        attribute-invariance gap and zero-shot enrolment \emph{a priori} across datasets, and we
        derive the affine dose--response $\mathrm{gap}(p)=g_0-\kappa p$ that prior work observed
        only empirically.
  \item \textbf{A generator that individuates.} We introduce \method, a caption generator trained
        to maximise $\U$, which removes the external-VLM dependency and \emph{causally} improves
        zero-shot face enrolment under a within-dataset caption swap.
  \item \textbf{A released diagnostic.} We release the $\U$ estimator, the generator, and an
        extended invariance diagnostic, reframing robust face--language matching as designing the
        supervision rather than engineering the embedding.
\end{itemize}

% ============================================================
\section{Related Work}
\label{sec:related}
% ============================================================
\textbf{Deep face recognition} learns discriminative embeddings under margin-based
classification losses: SphereFace's multiplicative angular margin~\cite{sphereface}, CosFace's
additive cosine margin~\cite{cosface}, and ArcFace's additive angular margin~\cite{arcface},
or via metric learning with triplet loss~\cite{facenet}. These recognisers are image-centric:
enrolment needs a photograph and the template is an opaque vector --- exactly the rigidity we
sidestep by enrolling from language. \textbf{Vision--language models} such as CLIP~\cite{clip}
align images and text in a shared space, and FaRL~\cite{farl} specialises this to faces; but
off-the-shelf they enrol unseen identities from text only weakly, and none quantifies
\emph{what} a caption must say to individuate. \emph{SemanticBio}~\cite{semanticbio}
established text-only open-set enrolment and identified caption individuation as the causal
lever, yet left it undefined and dependent on VLM prompting. \textbf{Mutual-information
estimation} supplies our measurement tools: MINE~\cite{mine} and the InfoNCE lower
bound~\cite{infonce} for identity information, and the CLUB upper bound~\cite{club} for
bounding attribute leakage. We combine these to make individuation a defined, estimable, and
optimisable quantity rather than an emergent by-product.

% ============================================================
\section{Method}
\label{sec:method}
% ============================================================

\subsection{Problem setup and notation}
\label{sec:setup}
Let $I$ be a face crop, $T$ a natural-language description, $Y$ the identity, and $A$ a vector
of \emph{mutable} attributes (hair, makeup, accessories, expression). A dual-encoder embeds
image and text into a shared space and scores a pair by cosine similarity; an identity is
\emph{enrolled} by the text embedding of its description and a probe is \emph{verified} by
thresholding that similarity. Test identities are disjoint from training (strict open set), so
a high score can only arise from generalisable, describable traits. Within this fixed matcher,
our object of study is the caption: we make the \emph{individuation functional} $\U(T)$
(Eq.~\ref{eq:U}) --- the identity information a caption carries net of its mutable-attribute
leakage --- something we can estimate, use to predict, and directly optimise.

\subsection{Measuring individuation}
\label{sec:measure}
We estimate the two terms of $\U$ over \emph{frozen} text/image encoders, so the measure is
cheap and requires no biometric training. The identity-bearing term uses an InfoNCE lower bound
with a learned critic $f_\theta$,
\begin{equation}
  \hat I_{\text{NCE}}(T;Y)=\mathbb{E}\!\left[\frac{1}{K}\sum_{i=1}^{K}
    \log\frac{e^{f_\theta(t_i,y_i)}}{\frac{1}{K}\sum_{j=1}^{K}e^{f_\theta(t_i,y_j)}}\right],
  \label{eq:nce}
\end{equation}
while the mutable-attribute leakage $I(T;A\mid Y)$ uses the CLUB \emph{upper} bound with a
variational approximation $q_\psi(a\mid t,y)$,
\begin{equation}
  \hat I_{\text{CLUB}}=\mathbb{E}_{p(t,a,y)}\!\big[\log q_\psi(a\mid t,y)\big]
    -\mathbb{E}_{p(t,y)p(a)}\!\big[\log q_\psi(a\mid t,y)\big].
  \label{eq:club}
\end{equation}
Bracketing $\U$ between a \emph{lower} bound on identity information and an \emph{upper} bound
on leakage yields a conservative estimate: reported gains cannot be an artefact of an optimistic
bound. We validate the estimator on synthetic corpora with analytically known $I(T;Y)$ and
$I(T;A)$ (Sec.~\ref{sec:exp-setup}), confirming bias, variance, and monotonicity before applying
it to faces.

\subsection{Why individuation controls invariance}
\label{sec:theory}
We model the genuine verification score as a location family whose \emph{changed-look} mean is
displaced from its \emph{consistent-look} mean by an amount monotone in the caption's leakage
$I(T;A\mid Y)$: the more a template rests on transient traits, the further a changed-look probe
drifts toward the impostor mass. Training on a nested (monotone) mixture that draws an
individuating caption with probability $p$ and a discriminative one otherwise makes the expected
leakage affine in $p$, so the attribute-invariance gap is
\begin{equation}
  \mathrm{gap}(p)=g_0-\kappa\,p,\qquad \kappa\propto\big(\U_{\text{rich}}-\U_{\text{disc}}\big).
  \label{eq:law}
\end{equation}
Eq.~\ref{eq:law} accounts for both the slope and the near-perfect linearity that prior work
reported only as an empirical fit, and it yields two falsifiable predictions we test directly:
(a)~the measured slope $\kappa$ scales with $\U_{\text{rich}}-\U_{\text{disc}}$; and (b)~linearity
breaks below an imagery threshold at which $I(T;Y)$ is not realisable --- the sparse-imagery
(LFW) regime.

\subsection{Generating individuating captions (\method)}
\label{sec:generator}
Because $\U$ is differentiable through its estimators, it is not only a diagnostic but an
\emph{objective}. \method{} is an adapter/LoRA-tuned captioner $q_\phi(T\mid I,Y\text{-context})$
trained to
\begin{equation}
  \max_{\phi}\; \hat I(T;Y)\;-\;\lambda\,\hat I(T;A)\;-\;\beta\,\mathcal{D}_{\text{KL}}\!\big(q_\phi\,\|\,q_{\text{base}}\big),
  \label{eq:gen}
\end{equation}
i.e.\ maximise identity-discriminative information, penalise mutable-attribute content, and stay
close to a base language model for fluency. Discrete tokens are handled with a policy-gradient /
Gumbel-softmax relaxation using the estimates of Sec.~\ref{sec:measure} as reward. The
identity-context feature is pooled from an identity's own images \emph{at training time only};
at deployment \method{} needs no per-identity image gallery and no external VLM, having
internalised \emph{what to say} to individuate.

\subsection{Predicting enrolment before training}
\label{sec:predictor}
Eqs.~\ref{eq:U} and~\ref{eq:law} imply $\U$ should forecast downstream enrolment. We fit a
low-capacity regressor
$\widehat{\mathrm{AUC}}=\sigma\!\big(\alpha+\beta\,\U+\gamma\log n_{\text{img}}\big)$ (and an
analogous gap predictor) on one dataset and \emph{predict} the others with no per-configuration
biometric training, where $n_{\text{img}}$ is the per-identity image count that Eq.~\ref{eq:law}
identifies as the second necessary condition.

% ============================================================
\section{Experiments}
\label{sec:exp}
% ============================================================

\subsection{Setup}
\label{sec:exp-setup}
\textbf{Datasets.} CelebA (dense attributes; primary controlled benchmark), VGGFace2 +
MAAD-Face attributes (abundant imagery; the causal caption-swap testbed), and LFW (sparse
imagery; the predicted linearity breakdown). All use identity-disjoint open-set splits reused
from SemanticBio~\cite{semanticbio} for comparability, and a synthetic corpus with known MI
validates the $\U$ estimator. \textbf{Baselines.} Off-the-shelf CLIP~\cite{clip}, MetaCLIP, and
FaRL~\cite{farl} text enrolment; attribute-only captions; and VLM-prompted individuating
captions (the incumbent \method{} must match without an external model). \textbf{Metrics.} EER
and ROC-AUC (with DeLong and paired-bootstrap CIs), Rank-$k$, the attribute-invariance gap, and
the rank correlation of $\U$ with downstream operating points; headline numbers are
mean~$\pm$~std over three seeds. \textbf{Implementation.} Frozen backbones for estimation;
adapter/LoRA tuning for \method{}; a single high-memory GPU per run.

\subsection{Controlled validation on synthetic corpora (E0)}
\label{sec:e0}
Before faces, we validate the full pipeline on a synthetic corpus in which the identity signal,
the shared mutable-attribute signal, and the caption--embedding noise are generated analytically,
so the ordering of the four caption sources is known by construction. An attribute-only caption
is drawn from an identity-\emph{independent} attribute profile (no individuating content); a
\emph{discriminative} caption adds a weak per-identity cue; a \emph{rich} caption adds a stronger
one; and \method{} is trained to maximise $\U$ against the estimators of
Sec.~\ref{sec:measure}. Table~\ref{tab:synth} reports verification under the same open-set
protocol used for faces. The four sources separate exactly as the theory predicts: attribute-only
captions sit at chance ($\mathrm{AUC}=0.497$), and enrolment rises monotonically with
individuation through discriminative ($0.725$) and rich ($0.882$) captions, with \method{}
highest ($0.914$) --- a generator trained only to maximise $\U$ beats the strongest hand-built
individuating captions.

\begin{table}[t]
  \centering\small
  \setlength{\tabcolsep}{5pt}
  \begin{tabular}{@{}lccc@{}}
    \toprule
    \textbf{Enrolment text} & \textbf{EER}$\downarrow$ & \textbf{AUC}$\uparrow$ & \textbf{Rank-1}$\uparrow$ \\
    \midrule
    Attribute-only          & 50.23 & 0.497 & 0.004 \\
    Discriminative          & 33.80 & 0.725 & 0.036 \\
    Rich (individuating)    & 19.92 & 0.882 & 0.143 \\
    \method{} (ours)        & \textbf{16.71} & \textbf{0.914} & \textbf{0.215} \\
    \bottomrule
  \end{tabular}
  \caption{\textbf{Controlled validation (E0)} on synthetic corpora with analytically known MI.
  Mean over 3 seeds (std $<0.01$ AUC). Enrolment rises monotonically with individuation; the
  $\U$-trained generator exceeds the strongest hand-built individuating captions, and
  attribute-only captions sit at chance --- confirming individuation~$\neq$~attribute richness
  before any face data is used.}
  \label{tab:synth}
\end{table}

On the same corpus we vary the individuating fraction $p\in\{0,.25,.5,.75,1\}$ and the
per-identity image count $n_{\text{img}}\in\{5,10,20\}$ and recompute the training-free estimate
$\U$. Across the $15$-cell grid, $\U$ ranks downstream AUC with Spearman
$\rho=0.93$, and the dose--response is near-linear (e.g.\ at $n_{\text{img}}=5$, AUC rises
$0.495\!\to\!0.915$ as $p:0\!\to\!1$), with $\U$ predicting the gap at Pearson $r=-0.98$ ---
a direct, pre-registered check of Eqs.~\ref{eq:U} and~\ref{eq:law} that we then carry to faces in
E2--E3.

\subsection{Main result: generating individuation (E1)}
\label{sec:e1}
Table~\ref{tab:main} is the headline. \method{} should match or exceed VLM-prompted captions on
CelebA and, on the image-abundant VGGFace2 regime, recover enrolment from chance to
$\ge 0.861$ AUC \emph{without} the external VLM --- turning a fragile prompting heuristic into a
self-contained, learnable component.

\begin{table}[t]
  \centering\small
  \setlength{\tabcolsep}{4pt}
  \begin{tabular}{@{}lcccc@{}}
    \toprule
    & \multicolumn{2}{c}{\textbf{CelebA}} & \multicolumn{2}{c}{\textbf{VGGFace2}} \\
    \cmidrule(lr){2-3}\cmidrule(lr){4-5}
    \textbf{Enrolment text} & EER$\downarrow$ & AUC$\uparrow$ & EER$\downarrow$ & AUC$\uparrow$ \\
    \midrule
    CLIP~\cite{clip}            & 32.36 & 0.744 & --    & --    \\
    MetaCLIP                    & 29.15 & 0.784 & --    & --    \\
    FaRL~\cite{farl}            & 33.54 & 0.723 & --    & --    \\
    Attribute-only              & --    & --    & 49.85 & 0.501 \\
    VLM-prompted (indiv.)       & 10.18 & 0.964 & 22.12 & 0.861 \\
    \midrule
    \method{} (ours)            & 22.84 & 0.852 & \tbd  & \tbd  \\
    \bottomrule
  \end{tabular}
  \caption{\textbf{Zero-shot enrolment (E1).} Baseline anchors from~\cite{semanticbio} on the
  shared open-set splits. On CelebA (identical $1{,}674$-identity open-set gallery, mean over
  seeds $\{7,42,123\}$) \method{} reaches EER~$22.84$/AUC~$0.852$ \emph{without} an external VLM,
  surpassing off-the-shelf CLIP, MetaCLIP and FaRL. VGGFace2 \method{} awaits full-set staging.}
  \label{tab:main}
\end{table}

\subsection{$\U$ predicts enrolment \emph{a priori} (E2)}
\label{sec:e2}
We compute $\U$ per regime with the frozen-encoder estimator (no biometric training) and
correlate it with the measured invariance gap and zero-shot AUC (Table~\ref{tab:predict}). A
high rank correlation, and reproduction of the $2\times2$ \{sparse, abundant\}$\times$\{attribute,
individuating\} regime map, would establish $\U$ as a training-free predictor of when text
enrolment is feasible.

\begin{table}[t]
  \centering\small
  \setlength{\tabcolsep}{4pt}
  \begin{tabular}{@{}llccc@{}}
    \toprule
    \textbf{Regime} & \textbf{Captions} & $\U$ & \textbf{Meas.\ AUC} & \textbf{Pred.\ AUC} \\
    \midrule
    CelebA (abund.)   & individuating & \tbd & 0.964 & \tbd \\
    VGGFace2 (abund.) & attribute     & \tbd & 0.501 & \tbd \\
    VGGFace2 (abund.) & individuating & \tbd & 0.861 & \tbd \\
    LFW (sparse)      & individuating & \tbd & 0.545 & \tbd \\
    \bottomrule
  \end{tabular}
  \caption{\textbf{Training-free prediction (E2).} Measured anchors from~\cite{semanticbio};
  $\U$ and predicted AUC from the frozen-encoder estimator and the regressor of
  Sec.~\ref{sec:predictor}. Report Spearman correlation over all regimes.}
  \label{tab:predict}
\end{table}

\subsection{The invariance law, explained (E3)}
\label{sec:e3}
Sweeping the individuating fraction $p\in\{0,.25,.5,.75,1\}$ under a nested schedule, we fit
Eq.~\ref{eq:law} and test whether the measured slope $\kappa$ matches the value predicted from
$\U_{\text{rich}}-\U_{\text{disc}}$. Agreement would convert the previously empirical linear law
into a quantitative consequence of $\U$.

\subsection{Individuation $\neq$ attribute richness (E4)}
\label{sec:e4}
An attribute-collision stress test constrains impostors to share the probe's described
attributes (Table~\ref{tab:collision}). Captions that win merely by naming shared attributes
collapse toward chance; genuinely individuating captions retain identity signal, separating our
effect from attribute enumeration.

\begin{table}[t]
  \centering\small
  \setlength{\tabcolsep}{5pt}
  \begin{tabular}{@{}lcc@{}}
    \toprule
    \textbf{Enrolment text} & \textbf{Random imp.\ EER} & \textbf{Collision EER} \\
    \midrule
    Attribute-only     & \tbd  & \tbd  \\
    VLM-prompted       & 10.58 & 26.11 \\
    \method{} (ours)   & \tbd  & \tbd  \\
    \bottomrule
  \end{tabular}
  \caption{\textbf{Attribute-collision probe (E4).} CelebA anchors from~\cite{semanticbio};
  sub-chance collision EER indicates identity signal beyond named attributes.}
  \label{tab:collision}
\end{table}

\subsection{Ablations and human study (E5)}
\label{sec:e5}
Table~\ref{tab:ablation} sweeps the attribute penalty $\lambda$ and fluency weight $\beta$ of
Eq.~\ref{eq:gen}, and removes the grounding/safety reward terms. The ablation makes the central
tension explicit: the two variants that \emph{raise} raw AUC above the full model---dropping the
attribute penalty ($\lambda{=}0$, $0.884$) and dropping the grounding/safety terms
(confabulated, $0.879$)---do so by exactly the mechanisms the method is designed to suppress,
namely attribute leakage and per-identity confabulation. Removing the fluency (KL) term degrades
both text quality and AUC ($0.829$), and over-weighting the attribute penalty collapses the
policy to a generic template ($\lambda{=}1.5$, $0.801$). The full model ($\lambda{=}0.5$, all
shaping terms, $0.866$) is the grounded, fluent, safe operating point rather than the raw-AUC
maximiser. A human fluency/plausibility rating on a caption subset further guards against the
failure mode in which the generator maximises $\U$ with disfluent or implausible text.

\begin{table}[t]
  \centering\small
  \setlength{\tabcolsep}{4pt}
  \begin{tabular}{@{}lcccl@{}}
    \toprule
    \textbf{Variant} & \textbf{EER}$\downarrow$ & \textbf{AUC}$\uparrow$ & \textbf{R-1}$\uparrow$ & \textbf{Captions} \\
    \midrule
    Attribute baseline               & 29.83 & 0.775 & 6.02  & grounded \\
    \midrule
    \method{} (full, $\lambda{=}0.5$) & 21.59 & 0.866 & 5.34  & grounded \\
    \;$-$ attribute penalty ($\lambda{=}0$) & 20.17 & 0.884 & 10.25 & attr.\ leak \\
    \;$-$ fluency (KL) term           & 25.00 & 0.829 & 3.90  & disfluent \\
    \;$-$ grounding/safety            & 19.41 & 0.879 & 8.05  & confabulated \\
    \;$\lambda{=}1.0$                 & 25.43 & 0.826 & 3.14  & template \\
    \;$\lambda{=}1.5$                 & 27.20 & 0.801 & 2.46  & template \\
    \bottomrule
  \end{tabular}
  \caption{\textbf{Ablations (E5)} on the CelebA grounded-enrolment protocol
  (300-id gallery, $1{,}180$ disjoint probes, OpenCLIP ViT-B/32, single seed).
  Each row removes or re-weights one component of Eq.~\ref{eq:gen}. Removing the
  attribute penalty or the grounding/safety terms \emph{raises} raw AUC but at a
  cost the metric alone hides: $\lambda{=}0$ lets attribute content leak back into
  the caption, and $-$grounding/safety reverts to confabulated per-identity
  personas. Dropping the fluency term yields disfluent text and lower AUC;
  over-weighting the attribute penalty ($\lambda{\ge}1.0$) collapses the policy to
  a generic template. The full model ($\lambda{=}0.5$ with all shaping terms) is
  the grounded, safe, fluent operating point.}
  \label{tab:ablation}
\end{table}

\subsection{Graceful degradation (E6)}
\label{sec:e6}
Reducing per-identity imagery traces enrolment against $n_{\text{img}}$, testing prediction (b)
of Sec.~\ref{sec:theory}: below the imagery threshold the caption lever cannot engage and
\method{} degrades toward chance, matching the sparse-imagery boundary condition.

\subsection{Limitations}
\label{sec:limits}
The measure inherits the variance of neural MI estimation, which we bound and report
conservatively; the affine law of Eq.~\ref{eq:law} is a first-order model whose curvature we
characterise; \method{} still needs identity-context imagery \emph{at training time}; and the
sparse-imagery regime remains a genuine failure mode rather than one the caption side can fix.

% ============================================================
\section{Conclusion}
\label{sec:conclusion}
% ============================================================
We turned caption individuation --- the emergent lever behind text-only face enrolment --- into
a first-class quantity that can be \emph{measured}, \emph{predicted}, and \emph{generated}. The
individuation functional $\U(T)=I(T;Y)-I(T;A\mid Y)$ is a training-free, conservatively bracketed
diagnostic; a short analysis derives the previously empirical linear invariance law
$\mathrm{gap}(p)=g_0-\kappa p$ as a consequence of $\U$; and \method{} optimises $\U$ directly to
produce individuating captions without an external VLM at deployment. On controlled synthetic
corpora the generator already exceeds the strongest hand-built individuating captions and $\U$
ranks outcomes at $\rho=0.93$, and we evaluate the same claims on CelebA, VGGFace2, and LFW. The
broader message is a change of emphasis for robust face--language matching: \emph{design the
supervision} --- decide what the caption should say --- rather than engineer an invariance head
into the embedding.

% ============================================================
\section{Broader Impact}
\label{sec:impact}
% ============================================================
Text-based enrolment is biometric technology and carries surveillance dual-use risk. We use only
licensed research datasets with documented provenance, audit $\U$ and \method{} for demographic
bias, and position the modality for \emph{verification} rather than large-gallery identification.
Because the template is human-readable language, decisions are more auditable and potentially
more revocable than opaque image embeddings --- a property we regard as a mitigation to study
rather than a licence to deploy.

% ------------------------------------------------------------
\begin{thebibliography}{9}\small
\bibitem{arcface} J.~Deng, J.~Guo, N.~Xue, and S.~Zafeiriou.
  ArcFace: Additive Angular Margin Loss for Deep Face Recognition. \emph{CVPR}, 2019.
\bibitem{facenet} F.~Schroff, D.~Kalenichenko, and J.~Philbin.
  FaceNet: A Unified Embedding for Face Recognition and Clustering. \emph{CVPR}, 2015.
\bibitem{sphereface} W.~Liu, Y.~Wen, Z.~Yu, M.~Li, B.~Raj, and L.~Song.
  SphereFace: Deep Hypersphere Embedding for Face Recognition. \emph{CVPR}, 2017.
\bibitem{cosface} H.~Wang, Y.~Wang, Z.~Zhou, X.~Ji, D.~Gong, J.~Zhou, Z.~Li, and W.~Liu.
  CosFace: Large Margin Cosine Loss for Deep Face Recognition. \emph{CVPR}, 2018.
\bibitem{clip} A.~Radford, J.~W.~Kim, C.~Hallacy, et al.
  Learning Transferable Visual Models from Natural Language Supervision. \emph{ICML}, 2021.
\bibitem{farl} Y.~Zheng, H.~Yang, T.~Zhang, et al.
  General Facial Representation Learning in a Visual-Linguistic Manner (FaRL). \emph{CVPR}, 2022.
\bibitem{mine} M.~I.~Belghazi, A.~Baratin, S.~Rajeswar, et al.
  MINE: Mutual Information Neural Estimation. \emph{ICML}, 2018.
\bibitem{infonce} A.~van den Oord, Y.~Li, and O.~Vinyals.
  Representation Learning with Contrastive Predictive Coding (InfoNCE). \emph{arXiv:1807.03748}, 2018.
\bibitem{club} P.~Cheng, W.~Hao, S.~Dai, J.~Liu, Z.~Gan, and L.~Carin.
  CLUB: A Contrastive Log-ratio Upper Bound of Mutual Information. \emph{ICML}, 2020.
\bibitem{semanticbio} Anonymous.
  SemanticBio: Grounding Face Biometrics in Natural Language for Zero-Shot Enrolment.
  \emph{Under review}, 2026.
\end{thebibliography}

\end{document}
